%% This document created by Scientific Word (R) Version 2.5
%% Starting shell: mitthesi


\documentclass[10pt,thmsa]{article}
%%%%%%%%%%%%%%%%%%%%%%%%%%%%%%%%%%%%%%%%%%%%%%%%%%%%%%%%%%%%%%%%%%%%%%%%%%%%%%%%%%%%%%%%%%%%%%%%%%%%%%%%%%%%%%%%%%%%%%%%%%%%
\usepackage{GECCO_proceedings}
\usepackage{natbib}

%TCIDATA{Created=Thu Mar 23 16:14:24 2000}
%TCIDATA{LastRevised=Wed Apr 04 12:04:36 2001}
%TCIDATA{Language=American English}

\input{
tcilatex}
\input{tcilatex}
\begin{document}

\title{The Mixture of Trees Factorized Distribution Algorithm}
\author{\textbf{Roberto Santana, Alberto Ochoa-Rodriguez, Marta R. Soto} \\
%EndAName
Center of Mathematics and Theoretical Physics. \\
ICIMAF. Calle 15, e/ C y D, Vedado CP 10400. C-Habana. Cuba\\
\{rsantana,ochoa,mrosa\}@cidet.icmf.inf.cu}
\maketitle

\begin{abstract}
This paper introduces a Factorized Distribution Algorithm based on a mixture
of trees distribution. The probabilistic model and the learning algorithm
used differs to previous uses of probabilistic modeling in the context of
Evolutionary Computation. Preliminary results show the algorithm is
competitive, and some times superior to other Factorized Distribution
Algorithms. We also illustrate how particular features of the search space
can be employed during the search by conveniently selecting the mixture of
trees parameters.
\end{abstract}

\medskip \medskip \medskip \medskip

\section{INTRODUCTION}

Factorized Distribution Algorithms (FDAs) \citep{Muhlenbein_et_al:1999} are
population based search methods that combine results from Graphical Models
and Evolutionary Computation research, and are considered as a tractable
subclass of Estimation Distribution Algorithms \citep{Muhlenbein:1996}. In
order to optimize a given function they begin by generating an initial
random population of points which are evaluated using the objective
function. Some of the points are selected based on their values, and a
factorized probabilistic model of their underlying distribution is
constructed. This probabilistic model is used to sample the points that will
be part of the next population.

A number of FDAs that use probabilistic models based on dependency trees
have shown up in the literature. In \citep{Baluja_and_Davies:1997} a tree
factorization computed using the Chow and Liu algorithm %
\citep{Chow_and_Liu:1968} is employed to approximate the underlying
distribution of the selected points. The Bivariate Marginal Distribution
Algorithm (BMDA) \citep{Pelikan_and_Muehlenbein:1999} uses a forest instead
of a tree based factorization. Only second order statistics are computed.

The Polytree Approximation Distribution Algorithm (PADA) has been designed
to deal with the class of single connected Bayesian Networks (BNs) %
\citep{Soto_et_al:1999}. A revised version of this algorithm has been
presented in \citep{Ochoa_et_al:2000} where the previous introduced
algorithms based on trees and forest distributions are covered.

Modeling by finite mixture of distributions \citep{Everitt_and_Hand:1981}
concerns modeling a statistical distribution by a mixture (or weighted sum)
of other distributions. Recently, the research on mixture models has begun
to receive a particular attention by the EDAs community. In %
\citep{Thierens_and_Bosman:2001} the authors use a mixture of Gaussian
probabilistic density functions for the solution of continuous
multi-objective functions. In \citep{Penha_et_al:2001} authors employ
mixture models as the basis for data clustering in multimodal continuous and
discrete function optimization via EDAs. In the case of discrete
optimization a framework for learning mixture of BNs that share the same
structure is presented.

Mixture of distributions have been used also in the framework of
evolutionary optimization that use flexible probability estimators %
\citep{Gallagher_et_al:1999} (an approach very close to FDAs) where the
adaptive mixture model of Priebe is used to estimate probability
distributions in an evolutionary optimization context.

\section{TREES AND MIXTURE OF TREES MODELS}

Mixtures of trees belong to the class of finite mixture distributions. They
were introduced in \citep{Meila:1999} and are the core of the optimization
algorithm we present in this paper. Now the probabilistic models are
formally introduced. We will utilize the same notation used in %
\citep{Meila:1999}.

Let $V$ denote the set of variables of our problem. According to the
graphical model paradigm, each variable is viewed as a vertex of an
(undirected) graph $G=(V,E)$ which is called a tree if it has no cycles. Now
we define a probability distribution $T$ that is conformal with a tree.

\begin{equation}
T(x)=\prod_{v\in V}T_{v|pa(v)}(x_v|x_{pa(v)})  \label{T-Tree}
\end{equation}

The distribution $T$ itself will be called a tree when no confusion is
possible. The graph $(V,E)$ represents the structure of the distribution $T$.

A mixture of trees is defined to be a distribution of the form:

\begin{equation}
Q(x)=\sum\limits_{k=1}^m\lambda _kT^k(x)  \label{Q-Mixture}
\end{equation}

with $\lambda _k\geq 0$, $k=1,..,m$, $\sum_{k=1}^m\lambda _k=1$.

The tree distributions are the mixture components, and the $\lambda _k$ are
called mixture coefficients. A mixture of trees can be viewed as containing
an unobserved choice variable $z$, which takes values $k\in $ $\{1,...,m\}$
with probability $\lambda _k$. Conditioned on the value of $z$ the
distribution of the visible variables $V$ is a tree. The $m$ trees may have
different structures and different parameters.

\section{THE LEARNING ALGORITHM IN THE CONTEXT OF THE MT-FDA}

We present first an algorithm for fitting a mixture of trees to an observed
data set in the Maximum Likelihood paradigm via the Expectation-Maximization
(EM) algorithm \citep{Dempster_et_al:1977}. The mixtures of trees learning
algorithm constitutes a building block for the FDA presented here, this
algorithm can be found in \citep{Meila:1999} where it was introduced. We
resume some of its features:

The learning problem is: Given a set of observations $D=%
\{x^{_1},x^{_2},,x^{_N}\}$, we are required to find the mixture of trees $Q$
that satisfies

\begin{equation}
Q=\arg \max_Q\sum\limits_{i=1}^N\log Q(x^i)  \label{argmax}
\end{equation}

The EM algorithm introduces a likelihood function called the $complete$ $%
log-likelihood$ which is the log-likelihood of both, the observed and the
unobserved data, given the current model estimate $M=\{m,T^k,\lambda
_k,k=1,...m\}$

\begin{equation}
lc(x^{1...N},z^{1...N}|M)=\sum_{i=1}^N\sum_{k=1}^m\delta _{k,z^i}(\log
\lambda _k+\log T^k(x^i))  \label{likehood-extended}
\end{equation}

where $\delta _{k,z^i}$ is equal to one if $z^i$ is equal to the $k$th value
of the choice variable, and zero otherwise.

The idea underlying the EM algorithm is to compute and optimize the expected
value of $lc$. In the context of population based evolutionary optimization
our objective is to find a factorization of the probability distribution of
the selected set of points using a mixture of trees. We have called our
algorithm Mixture of Trees FDA (MT-FDA).

In figure 1 the pseudo-code of the MT-FDA is presented. The first population
is randomly generated by the algorithm. From the current population, a
subset of points is selected. A mixture of trees $Q$ that fits the selected
set is found using the mixture of trees learning algorithm that takes as
parameters the number of trees, and a schedule with the number of learning
steps in each generation. New points are generated by sampling from $Q$. The
best elitism scheme is used, where all the individuals selected in the
current population pass to the next one. Different replacing strategies can
be used to combine points from the current population and new generated
points in the next population.

\medskip

\begin{center}
$\mathbf{M}${$\mathbf{ixture}$ $\mathbf{of}$ $\mathbf{trees}$ $\mathbf{FDA}$ 
}$\mathbf{(MT-FDA)}$
\end{center}

\begin{itemize}
\item  \textbf{STEP 0:} Set $t\Leftarrow 0$. Generate $N\gg 0$ points
randomly.

\item  \textbf{STEP 1:} Select a set $S$ of $k$ $<N$ points according to a
selection method.

\item  \textbf{STEP 2: }Calculate a mixture of trees $Q$ that approximates $S
$ using the mixture of trees learning algorithm.

\item  \textbf{STEP 3:} Generate $N-k$ new points sampling from $Q$.

\item  \textbf{STEP 4:} Combine in the new population the $k$ selected
points with the $N-k$ new points. Set $t$ $\Leftarrow t+1$

\item  \textbf{STEP 5:} If the termination criteria are not met, go to STEP 1
\end{itemize}

\medskip

Figure 1: MT-FDA

\medskip

The determination of the extent of learning to be allowed is a sensitive
issue for the MT-FDA. When sampling from the learned model, the two
traditional goals of an efficient search, exploitation and exploration, have
to be accomplished. A model that best approximates the data can less likely
generate, during the sampling step, points that belong to unexplored areas
of the search space. So, it is advisable to stop the learning algorithm
before the improvement in the likelihood ceases. The number of trees also
influence the likelihood. The design of strategies that set the appropriate
schedule for the learning steps, and of criteria to determine the number of
trees for a given optimization problem, are topics where further research is
required.

We test two methods for selecting a convenient starting mixture of trees.
Both methods begin by initializing each component of the mixture using the
Chow and Liu's algorithm. They differ in the second step. The first method
makes perturbations to the structure of each tree. The parent of one of the
leaf nodes is replaced by another randomly selected node. By making this
change on a leaf node the procedures guarantees no cycle will be formed in
the tree after the perturbation.

The second method keeps intact the tree structure and makes a perturbation
to the initial probability values $P^k(x^i)$, in order to let the learning
algorithm change the structure when fitting the perturbed probabilities.
Recall that, independently of the selected root, Chow and Liu's algorithm
guarantees that all trees are equivalent in their representation of the
data. By applying a small perturbation on $P^k(x^i)$ values, identicalness
among the trees is broken. Both methods guarantee solutions superior to
random initialization.

Mixture distributions have a number of distinctive attributes that make them
particularly appealing for their use in the framework of FDAs. Maybe the
most important is the possibility of representing, condensed in just one
model, different patterns of interactions among the variables of the
problem. In Bayesian Networks the change in one variable's value can
determine changes only in the parameters of other variables, not in their
structural relation. In mixture of finite distributions the structure of
dependencies among a set of variables can change depending on the values of
the choice variable they depend on. This fact can be used in a flexible way
for incorporating diverse search strategies into the MT-FDA.

Now we present some alternative ways the choice variable can be incorporated
by an FDA whose model is a mixture of trees. The analysis is divided in two
scenarios, when the choice variable is hidden, and when it is known.

\subsection{The choice variable z is hidden.}

Two cases can be distinguished in this scenario. The first is when, although
the choice variable is unknown, it is in fact one (or a small subset) of the
variables of the problem. In this case a fitter approximation of the
selected points by means of a mixture of trees could be found by identifying
the set of variables. There exist methods for identifying these variables
from the analysis of data \citep{Meila:1999}. Although we do not explore
this trend in the present work, we hypothesize that this identification task
can be inserted in the evolutionary process.

The second case, and the one extensively considered in our experiments, is
when there is no previous information about the choice variable, but we try
to fit the target distribution with a mixture that has $m$ trees, based on
an unknown variable that could help to cluster the space of solutions. As
the existence of a mixture of trees based on such a kind of choice variable
is an assumption, there could be cases where the mixture of trees does not
lead to an appropriate factorization.

\subsection{The choice variable z is known}

Again we analyze two particular cases: When the choice variable is in fact a
variable, (or the subset of the variables) of the problem. And, when the
choice variable $z$ does not belong to the set of variables of the problem
but we can decide its values based on a predefined criterium. We mention two
cases:

a) $z$ is related to a genotype clustering of the set of points

b) $z$ is related to a phenotype clustering of the set of points.

It has been shown \citep{Pelikan_and_Goldberg:2000} that for certain type of
problems a genotype clustering of the population can contribute to an
efficient search. The mixture of trees model can be used not only to cluster
the space of solutions but also to exchange information among the clusters
during the evolution. When solutions could be classified in regions based on
their structural similarities, and a mapping between solutions that belong
to the same class and their fitness evaluation exists, it is expected that a
mixture of trees based on a choice variable describing the classes could
lead to a good approximation.

\section{EXPERIMENTS}

\medskip The experiments were designed to illustrate the behavior of the
proposed algorithm and compare its performance with other FDAs. First we
introduce the functions that were employed. Then we present a comparison
between the MT-FDA and a tree based FDA. Some experiments are presented to
illustrate the difference in the performance of the Bayesian FDAs and the
MT-FDA. Finally an example is shown on the convenience of using one of the
variables of the problem as the choice variable.

Let $u$ be the number of bits turned on in the string $x$. A function of
unitation is a function whose value depends only on the number of ones on an
input string. The values of the strings with the same number of ones are
equal. Deceptive functions are defined as a sum of more elementary deceptive
functions $f_k$ of $k$ variables.\medskip 
\begin{equation}
f(x)=\sum\limits_{j=1}^lf_k(s_j)\text{,}
\label{Deceptive additively decomposed function}
\end{equation}

\medskip

where $s_j$ are non-overlapping substrings of $x$ containing $k$ elements.

\medskip

Function OneMax:

\begin{equation}
OneMax(x)=\sum\limits_{i=1}^nx_i  \label{OneMax function}
\end{equation}

Function $f_{dec}^3$:

\begin{equation}
f_{dec}^3=\left\{ 
\begin{tabular}{ccc}
$0.9$ & $for$ & $u=0$ \\ 
$0.8$ & $for$ & $u=1$ \\ 
$0.0$ & $for$ & $u=2$ \\ 
$1.0$ & $for$ & $u=3$%
\end{tabular}
\right.  \label{Deceptive function of order 3}
\end{equation}

Function $f_{3deceptive}$:

\begin{equation}
f_{3deceptive}(X)=\sum\limits_{i=1}^{i=\frac
n3}f_{dec}^3(X_{3i-2},X_{3i-1},X_{3i})  \label{f3deceptive}
\end{equation}

Function $Isotorus$:\medskip 
\begin{equation}
\begin{tabular}{|l|l|l|l|l|l|l|}
\hline
$u$ & $0$ & $1$ & $2$ & $3$ & $4$ & $5$ \\ \hline
$IsoT_1$ & \multicolumn{1}{|r|}{$m$} & $0$ & $0$ & $0$ & $0$ & $m-1$ \\ 
\hline
$IsoT_2$ & $0$ & $0$ & $0$ & $0$ & $0$ & $m^2$ \\ \hline
\end{tabular}
\label{Isotorus table}
\end{equation}

\begin{eqnarray}
&&F_{IsoTorus}=  \label{Isotorus} \\
&&\sum\nolimits_{i=2}^nIsoT_2(x_{up},x_{left},x_i,x_{right},x_{down}) \\
&+&IsoT_1(x_{1-m+n},x_{1-m+n},x_1,x_2,x_{1+m})
\end{eqnarray}

where $x_{up}$, etc., are defined as the appropriate neighbors, wrapping
around.

Function $BigJump$:

\begin{equation}
BigJump=\left\{ 
\begin{tabular}{ccc}
$u$ & $for$ & $0\leq u\leq n-m$ \\ 
$0$ & $for$ & $n-m\leq u\leq m$ \\ 
$k\cdot n$ & $for$ & $u=m$%
\end{tabular}
\right.  \label{BigJump Function}
\end{equation}

Function $Nf_{dec}^3$:

\begin{equation}
Nf_{dec}^3=\left\{ 
\begin{tabular}{c}
$\sum\limits_{i=1}^{i=\frac{n-1}3}f_{dec}^3(X_{3i-1},X_{3i},X_{3i}+1),$ $if$ 
$x_1=1$ \\ 
$\sum\limits_{i=1}^{i=\frac{n-1}3}(1-f_{dec}^3(X_{3i-1},X_{3i},X_{3i}+1),$ $%
oth.$%
\end{tabular}
\right. \text{,}
\end{equation}

\subsection{Comparison between a tree based FDA and the MT-FDA}

\smallskip First we make a comparison between a tree based FDA and the
MT-FDA. The tree based FDA has the same pseudo-code as the MT-FDA presented
in section 4, but only one tree found using the Chow and Liu's algorithm is
employed. The main difference between Baluja's algorithm and our tree based
FDA is that we do not update the bivariate probabilities multiplying by a
decay factor, instead in every generation, bivariate probabilities are
calculated from the selected set. We study the difference in the behavior of
the tree based FDA and the MT-FDA for diverse truncation values and
functions.

\begin{table}[htbp]
\begin{center}
$
\begin{tabular}{|c|c|c|c|c|c|c|}
\hline
Function & $n$ & $\tau $ & \multicolumn{2}{|c}{Tree} & \multicolumn{2}{|c}{
MT-FDA} \\ \hline
&  &  & S & Gen. & S & Gen. \\ \hline
& $30$ & $0.05$ & $37$ & $4.92$ & $9$ & $6.33$ \\ \hline
$f_{3deceptive}$ & $30$ & $0.1$ & $60$ & $5.73$ & $63$ & $6.66$ \\ \hline
& $30$ & $0.15$ & $72$ & $6.84$ & $75$ & $7.8$ \\ \hline
& $30$ & $0.2$ & $65$ & $8.28$ & $75$ & $9.28$ \\ \hline
& $30$ & $0.25$ & $85$ & $8.8$ & $95$ & $9.56$ \\ \hline
$Isotorus$ & $36$ & $0.05$ & $80$ & $3.49$ & $75$ & $4.05$ \\ \hline
& $36$ & $0.1$ & $84$ & $4.4$ & $93$ & $4.9$ \\ \hline
& $36$ & $0.15$ & $85$ & $5.19$ & $92$ & $5.78$ \\ \hline
& $36$ & $0.2$ & $84$ & $6.14$ & $90$ & $6.5$ \\ \hline
& $36$ & $0.25$ & $71$ & $6.86$ & $91$ & $7.38$ \\ \hline
\end{tabular}
$%
\end{center}
\caption{Numerical results for MT-FDA and tree based FDA}
\end{table}

Common parameters for the experiments were: Population size = 1000, Best
elitism, Maxgen =20. Stopping criteria were: the maximum number of
generations, and a total homogeneity in the selected population (i.e. all
the individuals were identical).

The schedule for learning was the following: In the first population there
were 10 learning steps to learn the mixture, in the following generations 5
learning steps. Recall that for the MT-FDA the trees from which the learning
algorithm is started are also found using the Chow and Liu's algorithm. So
this initialization can be seen as the best we could achieve with just one
tree (what the tree based FDA actually does), but in the following learning
steps the likelihood of data given by the initial trees is improved.

In Table 1 are presented results for two different functions. $n$ is the
number of variables, $\tau $ the truncation parameter, $S$ the number of
times the optimum was found in $100$ runs and $Gen$. the average number of
generations needed to find the optimum. Initial experiments were made for $%
f_{3deceptive}$ with $\tau =0.05$. Results were not particularly heart
warming. In the table it can be appreciated that for this truncation value
the tree based FDA finds the optimum four more times than when MT-FDA is
used. However, for, $\tau =0.25$ the MT-FDA clearly surpasses the tree with
an impressive 95 percent of success. Similar results were achieved for the
Isotorus function.

\subsection{Comparison between Bayesian FDAs and the MT-FDA}

In the general class of BNs the conditional probability of a variable x can
depend on a subset of variables and not on only one like in the subclass of
trees. All these variables are called parents. The complexity of the network
is related to the maximum number of parents any variable x can have. BNs
learning algorithms allow to incorporate constraints related with a maximum
number of parents or the network complexity.

In \citep{Etxeberria_and_Larranhaga:1999} BNs were introduced in the
framework of Evolutionary Optimization to learn, in every generation, a
factorization of the selected points. In this paper we compare our algorithm
with other two Bayesian FDAs. The Bayesian Optimization Algorithm (BOA)
introduced in \citep{Pelikan_et_al:1999} and the Learning Factorized
Distribution Algorithm (LFDA) \citep{Muhlenbein_and_Mahnig:1999}. Both
algorithms use a Bayesian metric to measure the goodness of every Bayesian
structure found, and a search procedure to search in the space of possible
structures. In the BOA it is possible to set the maximum number of parents
in the learnt network. For the LFDA we will refer to results appeared in %
\citep{Muehlenbein_et_Mahnig:1999a}. Contrary to the BOA, the LFDA uses the
BIC score to find the Bayesian structure and it uses a parameter $\alpha $
that allows to control the network complexity.

\begin{table}[tbph]
\begin{center}
\begin{tabular}{|l|l|l|l|l|}
\hline
Function & n & FDA & $\alpha $/ Trees & S \\ \hline
$OneMax$ & $30$ & LFDA & $0.75$ & $80$ \\ \hline
& $30$ & LFDA & $0.5$ & $38$ \\ \hline
& $30$ & LFDA & $0.25$ & $2$ \\ \hline
& $30$ & Tree &  & $89$ \\ \hline
& $30$ & MT-FDA & $2$ & $77$ \\ \hline
& $30$ & MT-FDA & $5$ & $19$ \\ \hline
& $30$ & MT-FDA & $10$ & $0$ \\ \hline
$BigJump(30,3,1)$ & $30$ & LFDA & $0.75$ & $100$ \\ \hline
& $30$ & LFDA & $0.5$ & $96$ \\ \hline
& $30$ & LFDA & $0.25$ & $58$ \\ \hline
& $30$ & Tree &  & $99$ \\ \hline
& $30$ & MT-FDA & $2$ & $97$ \\ \hline
& $30$ & MT-FDA & $4$ & $91$ \\ \hline
& $30$ & MT-FDA & $6$ & $77$ \\ \hline
\end{tabular}
\end{center}
\caption{Numerical results for MT-FDA and the LFDA}
\end{table}

A number of experiments were conducted to compare the behavior of the MT-FDA
with the LFDA. In this case we used an elitism parameter of 1 (i.e. only the
best individual is passed to the next population), truncation selection and
Maxgen =20. Stopping criteria for MT-FDA were the same that in the previous
experiments. In Table 2 column 4 refers to the number of trees and the
network density for the MT-FDA and the LFDA respectively. Later, it is
explained how these values were chosen.

A number of interesting features can be noticed from the analysis of the
results shown in Table 2. For these functions the number of trees in the
MT-FDA and the density of the BN in the LFDA play a similar role. This
analogy is evident for function $BigJump$, when the number of trees or the
density of the BN are increased results are poorer. In the table, values
used for the number of trees were selected trying to emulate the values for
the density of the BN used in experiments published in %
\citep{Muehlenbein_et_Mahnig:2000b}, but we do not claim we have attained a
perfect correspondence between the number of trees shown in Table 2 and the $%
\alpha $ values shown for the same functions. The analogy is related only to
the progression of values for both parameters.

Results for function $Bigjump$ are an example too that when structural
learning is done, simpler models can be better than more complex ones.
Regarding differences between the MT-FDA and the LFDA, and beyond the
analogy between the number of trees and the network density, results were
very similar.

It has been acknowledged that Bayesian FDAs are very sensitive to the
selection method used \citep{Pelikan_et_al:2000d}. Thus, when we claim here
that one algorithm is superior to the other for the analyzed functions, the
assertion is only valid for the particular selection strategy employed. We
have specified the selection methods and parameters used in every comparison
between Bayesian FDAs and the MT-FDA.

\subsection{Comparison between BOA and the MT-FDA}

We also compare our algorithm with the Bayesian Optimization Algorithm (BOA)
introduced in \citep{Pelikan_et_al:1999}. BOA uses a Bayesian metric to
measure the goodness of every Bayesian structure found, and a search
procedure to search in the space of possible structures. The implementation
of BOA used in this paper rests on the BOA platform available at %
\citep{Pelikan_et_al:2000d}. It employs the Bayesian Dirichlet equivalent
(BDe) metric, a greedy search procedure, and decision graphs for making the
search more efficient.

Table 3 shows the experimental results of the comparison between the BOA and
the MT-FDA for 2 functions. The parameters of the algorithm were: Population
size = 600, Tournament selection with tournament size = 4, the worst 85
percent of the current population is replaced by the new generated
individuals. The BOA was set to stop when the optimum was found, or an early
convergence occurred (univariate marginals higher than 0.95). The stop
criteria for MT-FDA were the same as in previous experiments, but note that
in general this setting is very different from the one used in the previous
experiments.

\begin{table}[htbp]
\begin{center}
\begin{tabular}{|c|c|c|c|c|c|}
\hline
Function & n & FDA & P/ T & S & Gen. \\ \hline
$Isotorus$ & $36$ & BOA & $2$ & $86$ & $7.29$ \\ \hline
& $36$ & BOA & $4$ & $81$ & $7.23$ \\ \hline
& $36$ & BOA & $6$ & $78$ & $7.4$ \\ \hline
& $36$ & Tree &  & $72$ & $7.96$ \\ \hline
& $36$ & MT & $2$ & $87$ & $8.22$ \\ \hline
& $30$ & MT & $4$ & $89$ & $7.83$ \\ \hline
& $36$ & MT & $6$ & $93$ & $7.38$ \\ \hline
$f_{3deceptive}$ & $30$ & BOA & $2$ & $78$ & $7.47$ \\ \hline
& $30$ & BOA & $4$ & $77$ & $8.12$ \\ \hline
& $30$ & BOA & $6$ & $77$ & $8.49$ \\ \hline
& $30$ & Tree &  & $37$ & $10.46$ \\ \hline
& $30$ & MT & $2$ & $54$ & $11.26$ \\ \hline
& $30$ & MT & $4$ & $60$ & $11.3$ \\ \hline
& $30$ & MT & $6$ & $69$ & $10.87$ \\ \hline
\end{tabular}
\end{center}
\caption{Numerical results for MT-FDA and the BOA}
\end{table}

In the table the fourth column represents the maximum number of parents and
the number of trees for the BOA and the MT-FDA respectively. Column five
represents the number of times the algorithm has converged in 100 runs. In
the experiments for the $f_{3deceptive}$ function BOA clearly overperformed
the MT-FDA. For this function two issues are worth to point out. The poor
performance of the tree based FDA, and the observation that when the number
of trees is increased the gap in the behavior between the MT-FDA and the BOA
is reduced. Nevertheless, by further increasing the number of trees the
MT-FDA does not achieve better results than BOA.

This situation is completely reversed for function Isotorus. For this
function, as the number of trees is increased MT-FDA overperforms BOA, when
the number of trees is 6 this difference is evident. The general conclusion
from these experiments is that for certain kind of functions MT-FDA can
overperform Bayesian-FDAs. One question remains open: which are the criteria
that allow to decide whether it is more convenient to apply one FDA or
another? We hypothesize MT-FDA can be more suitable for functions with
general and scattered overlapping between the variables, that could be
''covered'' by a set of relatively dependent probabilistic models. This
seems to be the case for the Isotorus function.

\subsection{\protect\medskip The role of the choice variable in the MT-FDA.}

\begin{table}[htbp]
\begin{center}
\begin{tabular}{|l|l|l|l|l|l|l|}
\hline
Function & n & N & z & T & S & Gen. \\ \hline
$Nf_{dec}^3$ & $31$ & $2500$ & unknown & $0.15$ & $24$ & $8.96$ \\ \hline
& $46$ & $3000$ & unknown & $0.15$ & $21$ & $11.57$ \\ \hline
& $61$ & $3000$ & unknown & $0.25$ & $12$ & $19.8$ \\ \hline
& $31$ & $1500$ & $x_1$ & $0.15$ & $30$ & $3.1$ \\ \hline
& $46$ & $1500$ & $x_1$ & $0.15$ & $28$ & $6.25$ \\ \hline
& $61$ & $2500$ & $x_1$ & $0.25$ & $23$ & $10.48$ \\ \hline
\end{tabular}
\par
\end{center}
\caption{Numerical results for the $Nf_{dec}^3$ function.}
\end{table}

As it was briefly discussed in previous sections the use of the choice
variable by the MT-FDA allows different and flexible ways of conducting the
search. We conducted experiments to evaluate the convenience of using one of
the variables of the problem as the choice variable. These experiments are
very preliminary but help to illustrate how the parameters of the mixture
model can be conveniently used in the FDA framework. Table 4 shows the
results for function $Nf_{dec}^3$. This function has many global optimal
that are triggered by variable $x_1$. When $x_1=1$ there is just one optimum
(all variables are set to 1), when $x_1=0$ there $3^{\frac n3}$ optima (all
possible combinations of 2 ones in every partition). For every setting 30
experiments were run using truncation selection and a maximum of 25
generations.

In the table it is shown a comparison between two MT-FDA with mixtures
models composed of two trees. The mixture model of the first MT-FDA uses $%
x_1 $ as its choice variable, for the second the choice variable is unknown
as was the case in the previous experiments. When $x_1$ is treated as the
choice variable the individuals that fulfill ($x_1=0)$ are approximated by
one tree, and those that satisfy ($x_1=1)$ by the other. In this case there
is not learning of the tree structures, only the mixture coefficients are
calculated as $\lambda _0=\frac{N_0}N,\lambda _1=\frac{N_1}N$, where $N_0$
and $N_1$ are respectively the number of individuals that satisfy $(x_1=0)$
and $(x_1=1)$. Nevertheless such a mixture allows the MT-FDA to focus on the
region of the space of solutions defined by the $x_1$ value. Eventually, as
evolution advances, one of the coefficients becomes $0,$ and a tree based
FDA is run from then on.

In Table 4 the improvements achieved by considering $x_1$ as the choice
variable can be seen. This is a simple example of how the choice variable
can be incorporated also as a tool for influencing the exploration and
exploitation purposes of the search.

\section{CONCLUSIONS AND FURTHER WORK}

In this paper we have introduced a FDA based on mixture distributions. The
MT-FDA is different to other simple connected based FDAs, and to Bayesian
FDAs too. Contrary to Bayesian FDAs the MT model allows that the structure
of dependencies among a set of variables changes. Compared to simple
connected FDAs, in the MT-FDA, dependencies between the parents of a given
node can be represented.

The algorithm exhibits a number of characteristics that places it far apart
from Gallagher's algorithm \citep{Gallagher_et_al:1999}. It is appropriate
for problems with integer representation that can have many variables. The
performance of Gallagher's algorithm is more sensitive to the number of
variables. The type of mixture model both algorithms employ is different,
the adaptive mixture model of Priebe \citep{Priebe:1994} is used in
Gallagher's algorithm.

Compared to the Estimation of Mixture of Distribution Algorithm (EMDA)
introduced in \citep{Penha_et_al:2001} the algorithm exhibits also several
differences. MT-FDA is not focused on the optimization of multiobjective
functions although it can be employed in multiobjective optimization too.
Another difference is that our mixture model allows components with
different structures. The version of the EM algorithm we use in this paper
is different to the one presented in \citep{Penha_et_al:2001}.

In this paper we have shown that the MT-FDA can overperform the tree based
FDA for the functions considered. The algorithm is also better than Bayesian
FDAs for some functions like Isotorus. Another achievement of this paper is
to have presented different ways to influence the search by manipulating
parameters related to the mixture of trees.

The computational complexity of the MT-FDA is mainly given by the Estimation
and Maximization steps of the (EM) algorithm. In each iteration of the
mixture of trees learning algorithm the running time of the estimation step
is $\Theta (mnN)$, and for the maximization step is $\Theta (mnr_{\max }^2)$
where $r_{\max }^{}$ is the maximum cardinality of the variables. Both steps
are computationally expensive. We have not treated this question in the
present work. Currently we are investigating how to reduce the computational
burden associated to the EM algorithm in the context of the MT-FDA.
Promising lines of research and further work are:

\begin{enumerate}
\item  To design efficient learning schedules that help to diminish the
number of evaluations.

\item  To extend the use of mixtures to deal with more complex probabilistic
models.

\item  The use of MT-FDA as a method to implement parallel and distributed
population based search.
\end{enumerate}

\newpage \textbf{Acknowledgments}

We are grateful for constructive comments of anonymous reviewers. This work
was funded by the Cuban Ministry of Science, Technology and Environment,
under the project Low Cost Evolutionary Algorithms.

\bibliographystyle{aaai}
\bibliography{thesbib}

\end{document}
